\section{A residue formula with independent twists}\label{sec:residue}

The next step is a scalar formula for the vacuum contraction. We keep the twists independent while taking residues so that the level intersections are transverse. The collision analysis below proves that the proposed residues exhaust the global residue theorem.

Write \([u]=2\sin u\), shifts indicated by integers mean multiples of \(l=\pi/8\); all variables are taken modulo \(8l\) (binomial products below descend even when single brackets are antiperiodic). The variables lie in the complex torus $\mathbb C/(8l\mathbb Z)$, identified with $\mathbb C^*$ by $p=e^{2it}$. The regular configuration torus excludes every shifted root--root and root--site coincidence. For this calculation put
\[
 V(d)=-\frac1{[d-1][d+1]},\quad G(d)=-\frac{[d]^2[d-2][d+2]}{[d-1][d+1]},\quad
 S(d)=\frac{[d+1]}{[d-1]},\quad R(d)=\frac{[d-2][d+1]}{[d+2][d-1]}
\]
Here $V(d)$ and $R(d)$ denote scalar functions of a root difference, a convention local to the residue calculation. Take \(N\) sites (indexed below by \(b\)), split into lists \(X,Y\), with no two sites differing by a shift. In \(N\) root variables put
\[
 D_i=\prod_b S(t_i-b)\prod_{j\ne i}R(t_i-t_j),\qquad
 \Phi=\prod_{i,b}V(t_i-b)\prod_{i<j}G(t_i-t_j)\prod_{i,x\in X}S(t_i-x)^{-1}.
\]
\begin{lemma}[Independent-twist trace]\label{lem:residue}
For generic independent nonzero parameters \(K_i\),
\[
 \sum_{D_i=-K_i}\frac{\Phi}{\det(\partial_{t_j}\log D_i)}
 =(-1)^N\sum_{\mathcal L}
 \mathrm{Res}_{\mathcal L}\Phi\prod_{i\ {\rm left}}(D_i/K_i).
\]
The left sums only in the regular configuration torus (no coincident shift orbits, no coincidences with site orbits); a term on the right assigns all labels injectively to ordered frozen slots: length 0,1,2 at each site, \(x-1,x-2\) (called left) or \(y+1,y+2\), with second slot occupied only if first occupied. Scalar residues are successive ordinary one-variable residues in \(t_i\), earlier slot first.
\end{lemma} 

\begin{proof}
We first compute the residue at one allowed next slot. To show that these local
residues exhaust the answer, we put every shifted collision and every toric
end on a compact boundary. On that model the order along an arbitrary
boundary ray is an additive function of its scale and collision depths,
apart from one explicit contribution from each level equation. Its zero
cases determine the polar groups in the global residue theorem. The last
step sums the resulting recursion.

\subsection{The recursive residues}

For a partial labeled slot assignment $F$ respecting the slot order, let $I$
be the free labels and let $h=|I|$. Write $\Phi_F$ for the scalar coefficient
obtained from $\Phi$ by the prescribed residues at the slots in $F$, including
$D_i/K_i$ when the label being frozen is left. Write $D_i^F$ for $D_i$ after
substituting the frozen slots. The form at this stage is
\begin{equation}\label{eq:recursive-form}
 \omega_F=\Phi_F\prod_{i\in I}\frac{D_i^F}{D_i^F+K_i}
                 \bigwedge_{i\in I}dt_i.
\end{equation}
The wedge uses the inherited label order. At $F=\varnothing$ this is the
original form, and at $h=0$ it is the scalar $\Phi_F$, which ends the
recursion. For $h>0$ put $H_i=D_i^F/(D_i^F+K_i)$. When several left
multipliers occur, they may equivalently be inserted before the iterated
residues: after insertion of its own multiplier an earlier variable has a
simple pole, and every later multiplier is regular there for generic values
of the other variables.

The free pair factors remain $G$ in $\Phi_F$ and $R$ in $D_i^F$. The
one-body factors at a left site $x$, up to nonzero constants independent of
the free variables, are as follows. Here $d=t_i-x$ and $s$ is the number of
slots already filled at $x$.
\[
 \begin{array}{c|cc}
 s&v_s(d)\text{ in }\Phi_F&d_s(d)\text{ in }D_i^F\\
 0&-1/[d+1]^2&S(d)\\
 1&[d+3]^2/[d+2]^2&[d+1][d+2]/([d][d+3])\\
 2&-[d+4]^2&[d+2]/[d+4]
 \end{array}
\]
Indeed freezing a left label at the additive slot $\xi$ multiplies a spectator's bare factor by
\[
 \frac{G(t_i-\xi)}{R(t_i-\xi)}
 =-\frac{[t_i-\xi]^2[t_i-\xi+2]^2}{[t_i-\xi+1]^2},
\]
and its level function by $R(t_i-\xi)$. At a right site the two factors are
$v_s(-d)d_s(-d)$ and $d_s(-d)^{-1}$, respectively. At the next left slot,
$D_i^F$ has a simple zero and $H_i\sim D_i^F/K_i$; at the next right slot,
$D_i^F$ has a simple pole and $H_i\sim1$. In each case the resulting form
has a simple pole with exactly the scalar residue specified in the lemma.
This also computes simultaneous next-slot residues when at most one label
is frozen at each site and the remaining variables are generic.

\subsection{Resolving collisions and coordinate boundaries}

We need a boundary that records both relative rates of escape and collisions
among variables at one scale. All bracket products are rational in
$p_i=e^{2it_i}$; their finite zeros and poles lie on
$p_i=q^k p_j$ or $p_i=q^k e^{2ib}$, where $q=e^{2il}$ and $k$ is taken
modulo $8$. Thus the following boundary contains every divisor relevant to
the form.

\begin{claim}[A compact model adapted to the level equations]\label{cl:compactification}
At each recursive stage, the regular configuration torus admits a smooth compactification, obtained by proper modifications of the projective ratio model, with a simple-normal-crossing boundary containing all coordinate and shifted-collision divisors. Every $D_i^F$ extends to a holomorphic map to $\mathbb P^1$. For generic independent nonzero twists, any collection of level divisors $D_i^F=-K_i$ meets every boundary stratum transversely whenever it meets it at all.
\end{claim}

Here and in the order calculation we suppress the superscript $F$ on the
level functions. The construction has three parts.

\paragraph{Relative scales.}
Adjoin the fixed label $p_*=1$, and take the closure of the graph of all
ratios $[p_a:p_b]\in\mathbb P^1$ for $a,b\in I\cup\{*\}$. For each ordering
$\sigma(0),\ldots,\sigma(h)$ of these labels there is an affine chart with
coordinates
\[
 u_a=\frac{p_{\sigma(a)}}{p_{\sigma(a+1)}},\qquad 0\le a<h.
\]
For $a<b$, the ratio $p_{\sigma(a)}/p_{\sigma(b)}$ is
$u_a\cdots u_{b-1}$, and the opposite ratio is its projective reciprocal.
These charts are copies of $\mathbb A^h$. They cover the closure: for any
convergent sequence of torus points, order the moduli of its labels along a
subsequence; all successive ratios in that order are finite, and the
projective ratio identities give the displayed chart. The model is therefore
smooth and compact. The divisors $u_a=0$ are labeled by cuts of the ordered
chain. They divide the labels into blocks at a common scale, within which
all ratios are finite and nonzero. The zero scale-ratio coordinates are
transverse to the coordinates inside those blocks.

The cluster resolution is an instance of the configuration-space and wonderful-model strategy of Fulton--MacPherson and De Concini--Procesi~\cite{FultonMacPherson1994,DeConciniProcesi1995}; Li's Theorems~1.2--1.3 give a general formulation~\cite{Li2009Wonderful}. We describe the charts explicitly, since the shifted collisions, scale boundaries and orders of the present form must all be checked.

\paragraph{Shifted collision clusters.}
Within a common-scale block, divide by one representative coordinate.
Near a point where some labels agree up to powers of $q$, choose their
compatible offsets $k_i$, and use $q^{-k_i}p_i/p_{\rm rep}-1$ as relative
coordinates. If the block is anchored at a site, divide instead by that site
and include it as a fixed label with coordinate zero. Distinct site orbits
ensure that a cluster contains at most one site; a site can occur only in
the scale block containing $*$. In these coordinates the shifted equations
through the point are ordinary equalities of coordinates, including equality
to zero for an anchor. Different coincidence classes give independent
diagonal arrangements, independent also of the scale-ratio coordinates.

A connected cluster $C$ is a subset of free labels and possibly one site,
with a consistent offset for each label modulo $8$; without a site only
offset differences are specified. Its diagonal is imposed on the projective
ratios themselves, so it remains defined over a scale boundary. Blow up
the strict transforms of these diagonals of codimension at least two,
starting with contained diagonals before the diagonals containing them.
Codimension-one diagonals are retained as boundary divisors.

Here is the local model for these blowups. For a smooth center
$z_1=\cdots=z_c=0$, the blowup is the closure of the graph of
$[z_1:\cdots:z_c]$. In its chart with the $r$th projective coordinate equal
to one,
\[
 z_r=s,\qquad z_j=s\,u_j\quad(j\ne r).
\]
The construction is proper because that graph closure is closed in the
product with $\mathbb P^{c-1}$. Changes of transverse equations by an
invertible matrix give the transition maps. For a full diagonal, the $u_j$
describe the relative positions modulo translation and scale, with two
distinct positions normalized. At a point of the exceptional divisor,
the remaining equalities split into smaller diagonal arrangements among
the labels with the same $u$-value; their equations do not involve $s$.
Away from the full diagonal they already split in this way. Induction on
the number of labels proves that every subsequent center is smooth and
that the resulting boundary has simple normal crossings. In particular,
clusters through one point are disjoint or nested, with compatible offsets;
the site is included as a label in this incidence rule.

We will use two consequences of this local model. A cluster of $n$ free
labels has codimension $n-1$ if unanchored and $n$ if anchored. The factor
of the pulled-back ordinary volume form in the displayed chart is
$s^{c-1}ds$, or $s^c\,ds/s$ in logarithmic coordinates. Thus its
logarithmic discrepancy is that codimension. Also, the valuation of a
binomial along the original cluster divisor is one exactly when its
equality is imposed by that diagonal, and zero otherwise. These statements
are read at a generic point of the cluster divisor, away from other
boundary components, and remain its valuations in all later charts.

\paragraph{Making the level maps regular.}
After the cluster blowups, every binomial in $D_i$ is a monomial in the
boundary parameters times a unit. In a chart with boundary parameters
$z_\alpha$ we may therefore write
\[
 D_i=u_i(z,y)\prod_\alpha z_\alpha^{c_{i\alpha}},
 \qquad u_i\ne0.
\]
The cone $\mathbb R_{\ge0}^{\{\alpha\}}$ records nonnegative combinations
of the divisors in that chart. Subdivide it by the hyperplanes
$\sum_\alpha c_{i\alpha}v_\alpha=0$, for every $i$. Do this compatibly on
faces. A rational triangulation, followed by lattice star subdivisions,
gives a unimodular subdivision and can leave every smooth face that needed
no cut unchanged. For completeness, after triangulation a nonunimodular
simplicial cone has a nonzero lattice point in the half-open parallelepiped
of its primitive rays. Subdividing along that point, also in adjacent cones
containing its minimal face, strictly lowers every changed determinant.
Choosing a cone of maximal determinant and repeating terminates. No such
point lies minimally in a smooth face, so the faces just specified remain
unchanged.

For a full-dimensional unimodular cone with columns
$a_{\alpha j}$, its chart is the monomial substitution
\[
 z_\alpha=\prod_j s_j^{a_{\alpha j}},
\]
with the other coordinates $y$ retained. On this chart the exponents of
each $D_i$ all have one weak sign. Hence either $D_i$ or $D_i^{-1}$ is
holomorphic, and $D_i$ defines a map to $\mathbb P^1$.

We record why these local substitutions give a proper modification.
On the overlap of two cones, the parameters for rays outside their common
face are invertible, and the transition maps are the corresponding Laurent
monomials. A point on that face is determined by its vanishing rays and by
the nonzero character monomials orthogonal to them. These gluings are
separated: modulo the common face, a linear functional separates the two
cones strictly, so a sequence of torus points cannot converge with a
vanishing ray outside that face in either chart and also converge in the
other. If all vanishing rays lie in the common face, the character
monomials identify the two limits. Density of the torus then gives
separatedness for the glued charts. Restricting to a coordinate face
downstairs gives exactly its subdivided fan with an extra torus factor;
unimodularity extends the face basis to a lattice basis. Consequently the
construction also glues between boundary charts: replacing any $z_\alpha$
by a unit times $z_\alpha$ changes each inverse-basis monomial only by a
unit.

To check properness, take a sequence whose image tends to a point of one
boundary chart and arrange $|z_\alpha|\le1$. Away from the boundary,
the vector $(-\log|z_\alpha|)_\alpha$ lies in a fixed subdivision cone
along a subsequence. In its unimodular basis the coefficients are
nonnegative, so all corresponding $|s_j|\le1$. A further subsequence
converges in that chart. For points over the boundary, approximate by
torus points in a compatible metric on the resulting smooth Hausdorff
manifold and apply the same argument. Thus the modification is proper.
The ratio model is compact and the cluster blowups are proper, so the
final model is compact as claimed.

Finally consider a subcollection of the level maps on a boundary stratum,
including the open stratum. Its critical values have measure zero. One
direct proof uses induction on the dimension of a coordinate chart:
on the locus of maximum rank this follows from the constant rank theorem;
the rest lies in the zero set of a nonzero maximal minor. Near every point
of that zero set, a suitable derivative of the minor has a smooth
hypersurface as its zero set, so countably many such lower-dimensional
patches cover it. A critical point remains critical after restriction,
and induction applies. Countable coordinate covers suffice. Excluding
these values for all strata, subcollections and recursive assignments
gives the asserted transversality for generic independent twists. This
proves Claim~\ref{cl:compactification}.

\subsection{Boundary orders and the surviving slots}

We now give the order calculation in a form that also applies to a ray
introduced by the last subdivision. A \emph{logarithmic order} is the
order of the coefficient when the form is written with $dz/z$ for each
boundary parameter $z$. Order zero permits a simple logarithmic pole;
strictly positive order removes that pole. The discrepancy computed above
is included in this convention.

\begin{claim}[Boundary orders]\label{cl:orders}
The pullback of $\omega_F$ has at worst logarithmic poles along the compactification boundary.
\end{claim}

\paragraph{The additive order and the level charges.}
Fix a final ray $\rho$. In its original boundary cone, write its
nonnegative cluster weights as $\lambda_C$ and write $w_i$ for the
valuation of $p_i$ contributed by its scale cuts, with $w_*=0$.
Only disjoint or nested clusters can have positive weight. The quantities
$|w_i|$ and $|w_i-w_j|$ are linear within this ordered scale cone.
The scale contribution of the bare form
$\eta_F=\Phi_F\bigwedge_{i\in I}dt_i$ is
\begin{equation}\label{eq:residue-scale-order}
 \mathfrak s_F(w)=h\sum_{i\in I}|w_i|-\sum_{i<j}|w_i-w_j|
       \ \ge\ \sum_{i\in I}|w_i|.
\end{equation}
Indeed every original site contributes $|w_i|$, and each of the $N-h$
frozen spectator factors removes $|w_i|$. A free pair contributes
$-|w_i-w_j|$, whereas $R,S$ and the logarithmic differentials have
scale order zero. The inequality is the triangle inequality summed over
pairs. In particular, a nonzero scale part has strictly positive order.

For a cluster $C$, let $m_k=m_k(C)$ count its free labels at offset
$k\in\mathbb Z/8\mathbb Z$, and put $n_C=\sum_km_k$ and
\begin{equation}\label{eq:residue-B}
 B(m)=\frac12\sum_{k\bmod8}(m_k-m_{k+1})^2
       +\sum_{k\bmod8}m_km_{k+2}.
\end{equation}
The pair order in $\Phi_F$ is
\[
 \sum_k m_k(m_k-1)+\sum_km_km_{k+2}-\sum_km_km_{k+1}
 =B(m)-n_C.
\]
Here the three terms come from equal offsets, offsets two apart, and
adjacent offsets in $G$. For a label at offset $k$, the pair contribution
to the valuation of $D_i$ is
\begin{equation}\label{eq:residue-pair-charge}
 c_{\rm pair}(k;m)=m_{k-2}+m_{k+1}-m_{k-1}-m_{k+2}.
\end{equation}
It sums to zero over all labels in the cluster, as also follows from
$R(d)R(-d)=1$.

For later use we record the site orders separately. Let
$\delta_{k,j}$ be one when $k=j$ modulo $8$ and zero otherwise. At a left
site of stage $s$, the bare order $\nu_s(k)=\operatorname{ord}_{d=k}v_s$
and the level charge $c_s(k)=\operatorname{ord}_{d=k}d_s$ are
\begin{equation}\label{eq:residue-site-charges}
\begin{array}{c|c|c}
s&\nu_s(k)&c_s(k)\\ \hline
0&-2\delta_{k,-1}&\delta_{k,-1}-\delta_{k,1}\\
1&2\delta_{k,-3}-2\delta_{k,-2}
 &\delta_{k,-1}+\delta_{k,-2}-\delta_{k,0}-\delta_{k,-3}\\
2&2\delta_{k,-4}&\delta_{k,-2}-\delta_{k,-4}.
\end{array}
\end{equation}
At a right site these are $\nu_s(-k)+c_s(-k)$ and $-c_s(-k)$.
There is no site order or charge at an unanchored cluster.

These data give an exact formula along $\rho$. In the next display
$\nu_C$ is the applicable bare site order just listed, or zero when
unanchored; $c_C$ is the applicable site charge. If label $i$ belongs
to $C$, its offset there is $k_i^C$. Then
\begin{align}
 \operatorname{ord}_{\log,\rho}\eta_F
 &=\mathfrak s_F(w)+\sum_C\lambda_C
   \left(B(m^C)-\mathbf1_{\{C\ {\rm unanchored}\}}
                  +\sum_km_k^C\nu_C(k)\right),\label{eq:residue-bare-order}\\
 e_i(\rho):=\operatorname{ord}_\rho D_i
 &=\sum_{C\ni i}\lambda_C
    \left(c_{\rm pair}(k_i^C;m^C)+c_C(k_i^C)\right),\label{eq:residue-level-order}\\
 \operatorname{ord}_{\log,\rho}\omega_F
 &=\operatorname{ord}_{\log,\rho}\eta_F
      +\sum_{i\in I}\max\{0,e_i(\rho)\}.\label{eq:residue-final-order}
\end{align}
The first line combines the pair order with the cluster discrepancy.
Logarithmic differentials transform by the determinant of the monomial
matrix, so their coefficient orders add with the weights $\lambda_C$.
The level functions have no scale valuation. Finally, at the generic
point of the ray divisor, $D_i=s^{e_i}u_i$ with $u_i$ a unit. Thus
$H_i$ has order $e_i$ if $e_i>0$ and zero if $e_i\le0$, away from its
separate level divisor. This proves the last line. In particular, the
maximum is taken after adding all depths in $e_i$.

\paragraph{A lower bound valid for every final ray.}
Let $A$ be the union of the free labels in the outermost positive-weight
clusters anchored at left sites. Use
\[
 \sum_i\max\{0,e_i\}\ \ge\ \sum_{i\in A}e_i.
\]
The outermost clusters are disjoint. Every positive-weight cluster is
contained in one of them or disjoint from it, by the cluster incidence
rule including the site labels. Therefore every pair charge cancels in
the sum on the right. Its only remaining terms are the site charges of
the left-anchored clusters inside those outermost clusters. Combining
them with the bare site orders in
\eqref{eq:residue-bare-order} gives
\begin{equation}\label{eq:residue-depth-bound}
 \operatorname{ord}_{\log,\rho}\omega_F
 \ \ge\ \mathfrak s_F(w)+\sum_C\lambda_C\,\ell(C),
\end{equation}
where the depth bounds are
\[
\begin{array}{c|l}
\text{cluster }C&\ell(C)\\\hline
\text{unanchored}&B-1\\
\text{left site, stage }0&B-m_{-1}-m_1\\
\text{left site, stage }1&B-m_{-2}-m_0+m_{-3}+m_{-1}\\
\text{left site, stage }2&B+m_{-2}+m_{-4}
\end{array}
\]
and right-site bounds are their reflections $k\mapsto-k$.

All these bounds are nonnegative, with a short explicit list of zero
cases. To see this, put $M=\max_km_k$ and
$\Delta=\frac12\sum_k(m_k-m_{k+1})^2$. If every $m_k$ is positive, the
distance-two sum is at least $\sum_km_k$ and strictly covers each
subtracted pair of counts. Otherwise
\[
 \Delta\ge\frac12\sum_k|m_k-m_{k+1}|\ge M.
\]
Equality requires one contiguous support interval, rising and falling
monotonically in steps of size at most one. For stage $0$, put
$a=m_{-1}$ and $b'=m_1$. If only one is positive, equality in
$B\ge M\ge a+b'$ requires that occupied offset to be a unit singleton
or part of a unit adjacent pair. If both are positive, use
$B\ge M+ab'\ge a+b'$; equality requires $a=b'=1$ and unit support
$\{-1,0,1\}$. Thus the stage-$0$ zero supports are exactly
\[
 \{-1\},\quad\{-2,-1\},\quad\{-1,0\},\quad
 \{1\},\quad\{0,1\},\quad\{1,2\},\quad\{-1,0,1\}.
\]
Every listed count is one. The stage-$1$ expression is the stage-$0$
expression translated to the offsets $-2,0$, plus
$m_{-3}+m_{-1}$. Its zero supports are therefore exactly
$\{-2\},\{0\},\{0,1\}$. Stage $2$ has no zero case for a nonempty
cluster. Finally $B=1$ has only a unit singleton or unit adjacent pair;
an unanchored cluster has at least two labels, so its zero case is a
unit adjacent pair. This proves Claim~\ref{cl:orders}.

\paragraph{Which zero-order rays can occur.}
For the total order in \eqref{eq:residue-final-order} to be zero,
\eqref{eq:residue-scale-order} and \eqref{eq:residue-depth-bound} require
$w=0$ and $\ell(C)=0$ at every positive depth. Equality in the bound
on the level terms also requires
\begin{equation}\label{eq:residue-equality-signs}
 e_i\ge0\quad(i\in A),\qquad e_i\le0\quad(i\notin A).
\end{equation}
These conditions, rather than the depth table alone, exclude most of its
zero cases.

For an outermost left-anchored cluster, the following table gives a label
whose charge from that depth is negative. The charge includes both
\eqref{eq:residue-pair-charge} and the site charge in
\eqref{eq:residue-site-charges}.
\[
\begin{array}{c|c|c|c}
\text{stage}&\text{outer support}&\text{tested offset}&\text{charge}\\ \hline
0&\{-1,0\}&0&-1\\
0&\{1\}&1&-1\\
0&\{0,1\}&1&-2\\
0&\{-1,0,1\}&1&-1\\
0&\{1,2\}&2&-1\\
1&\{0\}&0&-1\\
1&\{0,1\}&1&-1
\end{array}
\]
Every strict positive-weight subcluster must again be one of the zero supports
just listed if anchored, or a unit adjacent pair if unanchored. In the
stage-$0$ rows tested at offset $1$, the only possible anchored subclusters containing
that label are $\{1\}$ and $\{0,1\}$, with charges $-1$ and $-2$;
in the other rows there is no anchored subcluster containing the
label. Any unanchored subcluster containing a tested label is an adjacent pair
with that label at its right endpoint, and contributes $-1$. Hence every
eligible subcluster charge is nonpositive and the total $e_i$ of that label is negative,
contradicting \eqref{eq:residue-equality-signs}.

This leaves the next-slot singleton $\{-1\}$ at stage $0$, the next-slot
singleton $\{-2\}$ at stage $1$, and one extra stage-$0$ support
$\{-2,-1\}$. In that extra pair the charges at $-2,-1$ are respectively
$+1,0$. The only strict zero-order subcluster containing the label at $-2$
is the unanchored pair, where its charge is again $+1$; an anchored
$\{-1\}$ subcluster omits it. Thus the total charge of that extra label is
strictly positive. Right-anchored clusters give the reflected list and
opposite charges. An outermost unanchored zero support is an adjacent
pair with charges $+1,-1$ and has no strict subcluster; its positive
charge contradicts the second sign in
\eqref{eq:residue-equality-signs}. We have now classified every possible
boundary logarithmic pole, including rays combining several depths.

\begin{claim}[The contributing boundary intersections]\label{cl:surviving-slots}
Group each level pole by its free label. The boundary logarithmic poles can be assigned to those groups so that every intersection selecting one pole from each group consists of disjoint ordinary next-slot collisions and regular level intersections for the remaining labels. No additional collision or coordinate boundary contributes.
\end{claim}

Give group $i$ its level divisor $D_i=-K_i$. Every boundary logarithmic
pole has, in at least one outermost cluster, a label projecting to an
allowed next slot. Assign that component to the group of any such label;
for an extra stage-$0$ pair choose its next-slot label, not the extra one.
Claim~\ref{cl:compactification} implies that an intersection selecting
one component from every group is transverse and finite. If it selects
$r$ boundary components and $h-r$ levels, an extra boundary component
would leave a stratum of dimension less than $h-r$, while an extra level
would impose one more independent equation on the stratum of dimension
$h-r$. Transversality excludes both.

Suppose a selected component had an extra left pair, and let $i$ be its
label at $x-2$. Its total charge is positive, so the regular map $D_i$
is identically zero on that component. Its finite nonzero level cannot
be selected. Group $i$ would therefore have to select another boundary
component, on which $i$ projects to an allowed next slot. At their
intersection the two projections to the ratio model would agree, but
$x-2$ is not the stage-$0$ next slot $x-1$, and no other site's shift
orbit meets this one. This is impossible. At a right extra pair the
same argument uses $D_i=\infty$.

Every selected boundary ray therefore contains only next-slot
singletons. Such a cluster has no strict subcluster, and two such
clusters in one cone must be disjoint, including their site labels.
The cone they span is already smooth. Each level charge on it is
supported on its own label's slot ray, with one fixed sign, so none of
the level hyperplanes cuts that cone; the refinement kept it unchanged.
Thus its rays are the original codimension-one slot divisors, with at
most one per site. A point with only those boundary components cannot
project to any further collision or toric boundary: the inverse image
of such a divisor is supported on rays with that incidence. The model
is unchanged on this ordinary locus. The remaining levels are precisely
regular transverse solutions of the restricted system, which proves
Claim~\ref{cl:surviving-slots}.

\subsection{Completing the residue sum}

We apply the compact-manifold global residue theorem in the grouped-divisor
form of Griffiths~\cite[Section~III(a), Equation~(3.7)]{Griffiths1976Abel}.
Here compactness and the transverse finite intersections of the polar
groups were proved above. The sum of the residues in group order is zero.
One can see the sign directly: take residues along components of the
first $h-1$ groups, leaving meromorphic differentials on smooth compact
curves, and sum their ordinary residues. Poles at a second component
of one of those first groups occur twice with opposite signs, by
antisymmetry of the two residue orders. They cancel, leaving exactly
the sum over one component of every group. For $h=1$ this is the
ordinary residue theorem itself.

Let $T_F$ denote the all-level residue of $\omega_F$, summed over its
regular solutions; set $T_F=\Phi_F$ when $h=0$. At an all-level point,
the local residue is
$\Phi_F/\det(\partial_{t_j}\log D_i^F)$. By
Claim~\ref{cl:surviving-slots}, every other grouped intersection freezes
a nonempty subset $E$ of labels at distinct allowed next sites, then
takes the all-level residue of the remaining system. Conversely each
regular restricted solution occurs there. The local slot calculation
therefore gives the recursion
\begin{equation}\label{eq:residue-recursion}
 T_F=-\sum_{\substack{E\ne\varnothing\\
          E\ {\rm an\ allowed\ simultaneous\ next\ freeze}}}T_{F\cup E}.
\end{equation}
Permuting the selected groups and the corresponding variables changes
both orientations by the same sign, so no further sign occurs in this
formula. Vanishing local residues may be included.

Expand \eqref{eq:residue-recursion} to a full assignment $\mathcal L$.
Residues at slots of distinct sites commute: all interaction factors
between their disjoint shift orbits are regular, and the poles are
ordinary transverse coordinate poles. Along one site the first slot
must precede the second. Thus the scalar residue for $\mathcal L$ is
independent of its division into successive nonempty layers, while a
layering with $j$ layers has sign $(-1)^j$.

For a fixed assignment with $c$ two-slot chains and $s$ single slots,
$N=2c+s$. Let $a_j$ count its surjective layerings into $j$ ordered
layers, strictly increasing along each chain. For nonnegative integers
$p$, choosing the used layers gives the polynomial identity
\[
 \binom{p}{2}^{c}p^s=\sum_j a_j\binom{p}{j}.
\]
For the empty assignment the convention is $a_0=1$.
Evaluating at $p=-1$ yields
$\sum_j(-1)^j a_j=(-1)^s=(-1)^N$. Consequently every full assignment
has the coefficient in the lemma. The $h=0$ convention supplies the
terminal scalar residues, and the proof is complete.

\end{proof}
